% Compile with XeLaTeX and ko.TeX.
\documentclass[11pt]{article}
\usepackage[margin=25mm]{geometry}
\usepackage{amsmath}
\usepackage{kotex}
\title{RMFS 경로 선택과 스냅샷 선행 탐색}
\date{}
\begin{document}
\maketitle

\section{창고와 목표}
랙 클러스터 10개에 2열 6행씩 총 120개 랙을 배치하고 작업대 10개를 둔다. 로봇은 격자 경로에서 직각 회전을 한다. Makespan은 시뮬레이션에서 전체 작업이 끝난 시각이고, 실행 시간은 의사결정 계산에 걸린 실제 시간이다.

\section{DEVS 교차로 실험}
시뮬레이션 5초에 SOCS 스냅샷을 저장하고 분기마다 복구한다. 거리 기준 A* 최단경로에서 대기하는 Wait, 남북 우회, B 경로 회피, 양보점 세 곳을 비교한다. 서로 다른 유형 네 가지를 먼저 평가하고, 교통이 없는 조건의 최소 이동 시간으로 현재 최선을 갱신할 수 있는 후보를 추가 평가한다. 랙이 선분을 벗어나는 시각, 이동 완료, 예정된 진입을 사건으로 예약한다. 탐색에서는 연속 샘플링과 애니메이션을 끄고, 검증 재생에서는 0.01초 간격으로 랙 겹침을 검사할 수 있다.

\section{교차로 결과}
로봇 10대에서 진입 지연이 없으면 북쪽 우회는 58.30초, Wait는 59.76초에 완료했다. 진입을 0.25초 늦추면 단일 대안과 선행 탐색 모두 Wait를 선택했고 makespan은 59.76초였다. 실제 실행 시간은 단일 대안 0.316초, 탐색 2.082초였다. 7개 중 6개 분기를 실행해 완료 5개, dead-end 1개, 반복 상태 감지 0개, 시간 초과 0개를 기록했다. 분기 검사 시간은 총 2.070초, 평균 345 ms, 최대 404 ms였다. Dead-end는 STALLED로 집계했다. 정체 분기에서는 목적지에 정차한 랙이 교차 차량의 통로를 막았다.

\section{100주문 예약 실험}
시드 42로 생성한 주문 100개를 로봇 10대가 처리한다. 주문마다 홈에서 랙을 픽업해 작업대로 운반하고, 랙을 원위치에 반송한 뒤 홈으로 복귀한다. 교통에 따른 대기는 출발 전 홈에서 한다. 각 주문에서 기본 A* 최단경로와 선분 하나를 제외해 이동 구간 하나를 변경한 대안들을 같은 시공간 예약 상태에서 평가한다. 예상 홈 복귀가 가장 빠른 후보를 선택한다. 이 실험의 후보 평가는 예약 계산으로 수행하고, 선택한 스케줄은 DEVS로 재실행한다.

\section{100주문 결과}
후보는 총 701개로 주문당 평균 7.01개, 최소 5개, 최대 10개였다. 기준선은 주문당 단일 A* 계획이다. 기본 경로를 유지한 주문은 15개, 같은 거리의 대안을 선택한 주문은 78개, 더 긴 대안을 선택한 주문은 7개였다. 동거리 대안의 구간별 건수는 홈에서 랙까지 13개, 작업대까지 27개, 랙 반송 26개, 홈 복귀 12개였다. 더 긴 대안은 작업대 이동 2개와 랙 반송 5개였다. Makespan은 3433.63초에서 3111.06초로 9.4\% 감소했다. 계획 계산 시간은 0.803초에서 4.805초로 늘었고, 후보 검사 시간은 평균 6.3 ms, 최대 22.5 ms였다. DEVS 재실행에서 주문 100개를 완료했고 검사한 사건 시각에서 충돌은 검출되지 않았다.

\section{한계와 후속 실험}
계산 시간은 측정 환경에 따라 달라진다. 교차로 탐색은 한 시점에서 분기한다. 운행 중 반복 의사결정, 결정 마감 하에서의 상위 지연 시간, 일반적인 교착 방지 여부는 후속 검증 대상이다. 샘플 사이의 충돌 여부는 연속 궤적 검사로 확인해야 한다. 기하학적 공유와 시간상 경합의 영향을 구분하려면 후보별 교차점과 공유 선분 수, 점유 시간 중첩, 홈 대기, 회전, 추가 주행 거리를 기록한다.

\section{재현 명령}
{\small
\begin{verbatim}
python -m examples.RMFS.warehouse_crossings --robots 10 --stats
python -m examples.RMFS.warehouse_crossings --robots 10 `
  --crossing-delay 0.25 --stats
python -m examples.RMFS.fleet --robots 10 --orders 100 `
  --seed 42 --home-route-stats
\end{verbatim}
}

\section{운영 원칙과 다음 구간 계획}
로봇 10대와 주문 100개를 사용한다. 작업대는 실험 시작 전에 누적 출고 배정 건수가 가장 적은 곳으로 정하고, 동률이면 번호가 작은 곳을 선택한다. 모든 전략에 이 주문 목록과 작업대 배정을 동일하게 적용한다.

주문은 생성 순서대로 배정한다. 가용 로봇 중 출고 완료 건수가 가장 적은 로봇을 선택하고, 동률이면 번호가 작은 로봇을 선택한다. 로봇은 다른 주문에 배정되지 않은 랙 중 하나를 무작위로 선택한다. 시드 42를 바탕으로 주문별 무작위 랙 순서를 고정한다. 전략별로 당시 가용 랙이 다르면 실제 선택도 달라질 수 있다.

\begin{enumerate}
\item 가용 로봇에 주문을 배정하면 현재 위치에서 새 랙까지 계획한다.
\item 랙 아래에 도착하면 들어 올리기 직전에 배정 작업대가 비었는지 확인한다. 다른 로봇이 그 작업대를 향해 랙을 들어 올리기 시작했거나 작업대 점유가 계속되면 랙 아래에서 대기한다. 들어 올리기가 끝나면 작업대까지 계획한다.
\item 출고를 완료하면 작업대에서 랙의 원래 위치까지 계획한다.
\item 랙을 내려놓으면 대기 주문의 배차 대상이 된다. 대기 주문이 없으면 현재 위치에서 홈까지 계획한다.
\end{enumerate}

각 사건에서 다음 구간만 계획하고 이후 구간의 예비 예약은 두지 않는다. A*는 기본 최단경로 하나를 사용한다. A*k-X는 동거리 최단경로 후보를 다른 로봇의 활성 경로와 공유하는 노드 및 선분 수로 평가한다. A*k-L은 동거리 후보의 도착 시각을 스냅샷에서 비교하고, A*k-LD는 선분 제외로 얻은 우회 후보도 포함한다. 노드당 동률 이전 노드 한도는 $k=2$이고, 한 결정에서 검사하는 후보 상한은 $\ell=8$이다.

시공간 셀 예약은 주행 중의 점유를 검사한다. 로봇이 다음 구간을 확보할 때까지 도착 지점을 점유하며, 다음 구간이 당장 가능하지 않으면 그 자리에서 기다릴 수 있도록 모델링했다. 모든 활성 구간이 끝난 뒤에도 남은 로봇이 움직일 수 없으면 교착으로 기록한다. Makespan은 주문 100건 반송 및 최종 홈 복귀가 모두 끝난 시각으로 정의한다.

\begin{table}[htbp]
\centering
\caption{로봇 10대의 다음 구간 계획 실험}
\label{tab:next-leg-dispatch}
\begin{tabular}{rlrrrr}
\hline
번호 & 항목 & A* & A*k-X & A*k-L & A*k-LD\\
\hline
1 & Makespan (s) & NA & NA & NA & 2293.34\\
2 & 반송 완료 주문 & 53 & 42 & 97 & 100\\
\hline
\end{tabular}
\end{table}

A*k-LD는 100건을 반송하고 홈 복귀까지 완료했다. 나머지 세 전략은 진행할 구간이 없는 상태에서 종료됐다. 기본 A*는 종료 시 열 대 모두 다음 랙 접근 구간을 확보하지 못했다.

\begin{verbatim}
python -m examples.RMFS.leg_dispatch
\end{verbatim}

\end{document}
